\section{実験計画法の基礎とRCT}
\label{b17_design_RCT}

\subsection{授業内容}
\subsubsection{概要}
心理学研究では，変数間の因果関係を検証するために実験的アプローチが重要な役割を果たす。本コマでは，実験計画法の基礎的概念と，特にランダム化比較試験(Randomized Controlled Trial; RCT)の意義と方法について学ぶ。RCTは因果関係を検証するための強力な方法であり，その理論的背景と実践的手法を理解することは，心理学研究の質を高めるために不可欠である。また，実験における誤差の概念とその分布特性についても学び，統計的推論との関連を理解する。実験計画と統計的分析の橋渡しとなる基本的な考え方を身につけることを目指す。

\subsubsection{コマ主題細目}
\begin{description}
	\item[実験計画法の基本的考え方] 心理学における実験的アプローチの意義と限界について理解する。実験とは，操作した独立変数が従属変数に与える効果を測定することで因果関係を検証する方法である。相関研究と実験研究の違い，因果関係を主張するための条件(時間的先行性，共変関係，他の説明可能性の排除)について学ぶ。また，実験の内的妥当性と外的妥当性のバランス，実験における倫理的配慮についても理解する。
	      \begin{flushright}
		      $\to$ \textcite{Miura201705}のP.30--34，\textcite{Toyoda1992}のP.24-66.
	      \end{flushright}

	\item[ランダム化の意義と方法] ランダム化(無作為化)は実験研究の核心部分であり，偶然誤差や交絡要因の影響を統制するために不可欠な手法である。ランダム化比較試験(RCT)では，研究参加者を無作為に実験群と統制群に割り当てることで，群間の初期状態の等質性を確保する。これにより，介入効果以外の要因による影響を最小化し，観察される群間差を介入効果として解釈できるようになる。ランダム化の具体的な方法(単純ランダム化，ブロックランダム化，層別ランダム化など)とその利点・欠点について学ぶ。
	      \begin{flushright}
		      $\to$ \textcite{Yasui2019}を参照，RCTの基本的な考え方については\textcite{Toyoda201706}のP.5-26.
	      \end{flushright}

	\item[実験における誤差と正規分布] 実験データには必ず誤差(偶然変動)が含まれる。誤差は様々な要因の集積によって生じるが，中心極限定理により，多くの独立した小さな誤差要因が加算されると，その合計は正規分布に近似する。この性質が実験データの統計的分析の理論的基盤となっている。誤差分布としての正規分布の特性(平均0，対称性，確率的性質)を理解し，実験計画における誤差の扱い方と統制方法について学ぶ。また，誤差を分散という指標で評価することの意味と，それが実験効果の検出力にどのように影響するかについても理解する。
	      \begin{flushright}
		      $\to$ 誤差分布については\textcite{Naganuma2201611}のP.16--40，\textcite{Yamada2004}のP.80--85.
	      \end{flushright}

	\item[実験効果の評価方法] 実験効果を評価するためには，介入による変化(効果)と偶然による変動(誤差)を区別する必要がある。効果と誤差の分離は，線形モデルの枠組みで理解できる(個々の観測値 = 全体平均 + 介入効果 + 誤差)。効果の大きさは，群間差(実験群平均 - 統制群平均)で評価するが，その統計的有意性は誤差の大きさとの相対比較によって判断される。実験効果の評価における統計的検定の役割と，効果量(標準化された効果の指標)の重要性について理解する。
	      \begin{flushright}
		      $\to$ \textcite{kosugi2018}のP.60--67，効果量については\textcite{Matia201201}のP.43--46.
	      \end{flushright}

	\item[実験計画の種類と要素] 実験計画を構成する基本要素として，要因(factor)と水準(level)の概念を理解する。一要因計画と多要因計画の違い，群間計画(between-subjects design)と群内計画(within-subjects design)の特徴と選択基準について学ぶ。また，各種実験計画の表記法(例：「2$\times$3の混合計画」)についても理解する。次回以降の講義で詳細に学ぶ群間計画と群内計画の基本的な考え方について概観し，それぞれの利点と欠点を理解する。
	      \begin{flushright}
		      $\to$ \textcite{Toyoda201706}のP.27--44，実験計画の種類については\textcite{kosugi2018}のP.74--78.
	      \end{flushright}

\end{description}
\subsubsection{キーワード}
\begin{itemize}
	\item ランダム化比較試験(RCT)
	\item 実験群と統制群
	\item 偶然誤差と正規分布
	\item 要因と水準
	\item 群間計画と群内計画
\end{itemize}
\subsection{授業情報}
\paragraph{コマの展開方法}
講義
\paragraph{標準シラバスにおける位置づけ}
科目番号4；心理学研究法；1.心理学における実証的研究法(量的研究及び質的研究)；C 実証の手続き
\subsubsection{予習・復習課題}
\paragraph{予習}
実験計画法の基本的な考え方について，心理学研究法の教科書等で予習しておくこと。特に因果関係と相関関係の違い，実験研究と調査研究の違いについて理解しておくと講義の理解が深まる。また，前回学んだt検定の基本的な考え方を復習しておくと，実験計画と統計的分析の関連について理解しやすくなる。

\paragraph{復習}
授業で学んだ内容を定着させるために，以下の点について理解を深めておくこと：

\begin{enumerate}
\item ランダム化の意義と具体的な方法について説明できるようにする
\item 偶然誤差がなぜ正規分布に従うと考えられるのか，その理論的根拠を理解する
\item 実験効果と誤差を分離する考え方を理解し，線形モデルの枠組みで説明できるようにする
\item 様々な実験計画の種類とその特徴を理解し，具体的な研究例を挙げて説明できるようにする
\item 実験研究の内的妥当性と外的妥当性のバランスについて考察する
\end{enumerate}

また，次回の授業で学ぶ群間計画(Between-subjects design)の準備として，\textcite{Toyoda201706}のP.27--44や\textcite{kosugi2018}のP.74--78を読んでおくとよい。